\chapter{/boot：从 bootpack 视图到可读写 FAT32}
\label{chap:bootfat32}

本章说明 \file{/boot} 的两代实现：改动前的 bootfs（bootpack 交接视图）和现在的 FAT32 文件系统（xv6 风格的 bio + fat32 分层），以及它们和 SPL、XR829 固件加载、tftp 的关系。根文件系统 minifs、VFS 挂载表、devfs 的整体结构见第~\ref{chap:rootfs} 章。

一句话总结：

\begin{itemize}
  \item \textbf{改动前}：读 FAT32 的是 SPL，内核里的 \file{/boot} 只是把 SPL 预先搬进 DDR 的三份固件包装成文件，只读、只认三个写死的路径、不能列目录。
  \item \textbf{改动后}：内核自己读写 SD 卡上的 FAT32 分区，\file{/boot} 是真正的文件系统，可以按文件名查找、列目录、创建和写入文件；tftp 下载的文件直接落盘到卡上。老的 bootfs 只在 FAT32 挂载失败时作为后备。
\end{itemize}

\section{改动前：bootfs 只是 bootpack 的视图}

\subsection{加载流程}

\begin{txtcode}[basicstyle=\ttfamily\scriptsize]
miniBoot SPL  load_bootpack()
+-- 读 SD 128 KiB bootpack，装 kernel → 0x40200000
+-- bootpack header 拷到 0x401ff000
+-- load_fat32_firmware()                 SPL 里的只读 FAT32 解析器
    +-- fat32_mount()                     MBR 找 FAT32 分区，解析 BPB
    +-- fat32_load_root() × 3             只在根目录按名字找，沿簇链读入 DDR
    |   +-- boot_xr829.bin → 0x40100000
    |   +-- fw_xr829.bin   → 0x40110000
    |   +-- sdd_xr829.bin  → 0x40140000
    +-- bootpack_slot()                   往 header 补 3 个条目（offset=0 表示常驻镜像），重算 CRC

miniOS 内核  bootfs.c（63 行）挂在 /boot
+-- lookup("/boot/fw_xr829.bin")
|   +-- 路径写死映射 → t113_bootpack_find("fw_xr829")
|       → node.id = load_addr（DDR 物理地址当文件 id）
+-- read() = 从 node.id + offset 处 memcpy
\end{txtcode}

\subsection{本质}

bootfs 本质上是\textbf{bootpack 解析器}，挂到 \file{/boot} 只是借用 VFS 的接口，让 XR829 驱动可以用 \code{t113\_vfs\_open("/boot/fw\_xr829.bin")} 拿固件，而不用知道 bootloader 私有的地址常量。它不是文件系统：

\begin{table}[H]
\centering
\begin{tabular}{llp{6.6cm}}
\toprule
能力 & bootfs & 原因 \\
\midrule
按文件名查找 & $\times$ & \code{entry\_name()} 只认三个硬编码路径，映射到 bootpack 条目名 \\
列目录（\code{ls /boot}） & $\times$ & \code{readdir = 0} \\
读任意文件 & $\times$ & 只有 SPL 预加载的三份固件 \\
写入 / 创建 & $\times$ & 没有 \code{write}、\code{create} \\
看到运行期卡上的变化 & $\times$ & 数据是启动时的 DDR 快照，内核从不访问 FAT32 分区 \\
\bottomrule
\end{tabular}
\end{table}

所以当时 tftp 之类``把文件存进 \file{/boot}''的功能无从实现。

\section{改动后：分层结构}

照 xv6 的 \file{bio.c} / \file{fs.c} 分层：

\begin{txtcode}
用户态  open(O_CREAT|O_TRUNC) / read / write / readdir / tftp
   |
VFS     src/kernel/fs/vfs.c        挂载表最长前缀匹配，/boot → fat32_ops
   |
fat32   src/kernel/fs/fat32.c      路径 → 目录项，文件偏移 → 簇链；inode 表
   |
bio     src/kernel/fs/bio.c        512 B 扇区缓存，24 块 LRU，写直达
   |
SD      src/drivers/mmc/sunxi-sd.c t113_sd_read / t113_sd_write（CMD17/CMD24，轮询）
\end{txtcode}

\subsection{挂载顺序}

\file{src/kernel/init/main.c} 的 \code{kernel\_boot()}：

\begin{ccode}
t113_vfs_init();
/* The root device initialises the SD card that /boot shares. */
if (t113_minifs_mount_rootdev() < 0)        /* 先挂 /，顺带初始化 SD 卡 */
    kernel_platform_panic("cannot mount root filesystem");
if (t113_fat32_mount_vfs() < 0) {           /* 再挂 /boot（FAT32 读写） */
    t113_uart_puts("fat32: mount failed, using SPL handoff images\n");
    if (t113_bootfs_mount_vfs() < 0)        /* 后备：老的 bootpack 视图 */
        kernel_platform_panic("cannot mount boot filesystem at /boot");
}
if (t113_devfs_init() < 0) ...
\end{ccode}

与之前相比，\file{/} 改为先挂（rootdev 里会做 \code{t113\_sd\_init()}，\file{/boot} 和 \file{/} 共用同一张卡、同一个 SMHC0 控制器）。两者都在调度器启动前完成，此时互斥锁实际不生效，SD 访问是串行的。

\section{bio：扇区缓存（\file{bio.c}）}

\begin{table}[H]
\centering
\begin{tabular}{ll}
\toprule
项 & 值 \\
\midrule
块大小 \code{BSIZE} & 512 B（= SD 扇区，块号就是卡上的绝对 LBA） \\
缓冲数 \code{NBUF} & 24 \\
替换策略 & 双向链表 LRU，\code{brelse} 后放到表头 \\
写策略 & \textbf{写直达}：\code{bwrite()} 立刻 \code{t113\_sd\_write()} \\
对齐 & \code{data[]} 4 字节对齐（SD 驱动用 CPU 搬 FIFO，按字访问） \\
\bottomrule
\end{tabular}
\end{table}

接口和 xv6 一致，多了一个 \code{bget\_blank()}：

\begin{table}[H]
\centering
\begin{tabular}{lp{9.2cm}}
\toprule
函数 & 作用 \\
\midrule
\code{bread(lba)} & 返回带引用的缓冲，未命中时从卡上读 \\
\code{bget\_blank(lba)} & 返回带引用的缓冲但\textbf{不读卡}，给``整块覆盖''的调用者（写满一个扇区、清零新簇）省一次读 \\
\code{bwrite(b)} & 写回卡上 \\
\code{brelse(b)} & 释放引用，缓冲留在缓存里 \\
\bottomrule
\end{tabular}
\end{table}

同一块的并发访问由上层串行化（fat32 每个操作都持有自己的锁），所以 bio 不像 xv6 那样给每个 buf 配 sleeplock。

因为是写直达，卡上的内容永远不比缓存旧，fat32 读大段数据时可以绕开缓存直接读卡（见下文``读路径''）。

\section{fat32（\file{fat32.c}）}

\subsection{挂载 \code{t113\_fat32\_mount\_vfs()}}

\begin{enumerate}
  \item \code{binit()}，\code{bread(0)} 读 MBR，验 \code{55 AA}，取\textbf{第一个}类型为 \code{0x0B/0x0C/0x1B/0x1C}（FAT32 CHS/LBA，含隐藏变体）的分区表项。
  \item 读分区首扇区（BPB），校验：\code{55 AA}、每扇区 512 字节、每簇扇区数非零、FAT 个数非零、根目录项数为 0（FAT32 特征）、每 FAT 扇区数和总扇区数非零。
  \item 算出布局：
\end{enumerate}

\begin{txtcode}
part_lba
+-- 保留区（含 BPB、FSInfo）         fat_lba  = part + BPB[14] 保留扇区数
+-- FAT × fats                       data_lba = fat_lba + fats × sectors_per_fat
+-- 数据区                           簇 n 的 LBA = data_lba + (n-2) × spc
\end{txtcode}

\begin{enumerate}
  \setcounter{enumi}{3}
  \item 簇数 = (总扇区 $-$ 元数据扇区) / spc，并用 FAT 表能描述的条目数封顶；根目录簇号必须合法。
  \item 挂到 \file{/boot}，打印 \code{fat32: /boot mounted rw, partition lba=… clusters=… cluster-bytes=…}。
\end{enumerate}

\subsection{内存中的 inode 表}

FAT32 没有 inode，文件的元数据（首簇、大小）存在目录项里。为了对齐 xv6 的 icache，fat32.c 维护一个 16 项的 \code{fat\_inode} 表：

\begin{ccode}
struct fat_inode {
    uint32_t ref;
    uint32_t dir;          /* 1 for a directory */
    uint32_t first;        /* first cluster, 0 for an empty file */
    uint32_t size;
    uint32_t ent_lba;      /* sector of the short entry, 0 for the root */
    uint32_t ent_off;      /* byte offset of that entry in the sector */
    uint32_t dirty;        /* first/size changed: rewrite the entry */
    uint32_t cur_index;    /* last cluster looked up: index in chain... */
    uint32_t cur_cluster;  /* ...and its number (0 = none) */
};
\end{ccode}

\begin{itemize}
  \item \textbf{身份} = 短目录项所在的 \code{(ent\_lba, ent\_off)}，同一个文件被打开两次会拿到同一个 inode（\code{iget} 按位置查找，引用计数 +1）。
  \item \code{vfs\_node.private\_data} 指向这个 inode，VFS 后续的 read/write/close 都通过它找到文件。
  \item \code{cur\_index / cur\_cluster} 缓存上一次查到的簇，顺序读写时沿簇链只往前走一步，不用每次从首簇开始数（簇链是单向链表）。
  \item \code{dirty}：写入导致首簇或大小变化时置位，\textbf{最后一次 close 时} \code{iupdate()} 把首簇（高 16 位 @20、低 16 位 @26）和大小（@28）写回目录项。
\end{itemize}

\subsection{路径解析}

\code{resolve(path)} 先用 \code{mount\_relative()} 去掉 \file{/boot} 前缀，然后逐个分量在目录簇链里用 \code{dir\_find()} 查找：

\begin{itemize}
  \item 目录遍历器 \code{dir\_next()} 同时拼接长文件名（LFN）片段，并用短名校验和（\code{short\_checksum()}）确认 LFN 和后面的短目录项属于同一个文件，否则退回用 8.3 短名。
  \item 名字比较不区分大小写；短名会按 NT 保留字节（byte 12）的大小写标志还原成小写。
  \item 支持任意层子目录；\code{..} 指向根目录时首簇为 0，按根目录簇处理。
  \item 根目录本身用 \code{lba = 0} 表示（它没有目录项）。
\end{itemize}

\subsection{读路径}

\begin{txtcode}
fat32_read(node, offset, dst, count)
+-- 循环
    +-- ibmap(ip, pos / 簇大小, alloc=0)   沿簇链找到所在簇
    +-- 扇区对齐 + 剩余 ≥ 512 + dst 4 字节对齐
    |   +-- t113_sd_read(lba, n, dst) 直接读进调用者缓冲（一次读完簇内剩余整扇区）
    +-- 否则 bread(lba) → copy 部分扇区 → brelse
\end{txtcode}

读大文件（固件、内核镜像）时绝大多数数据走直读，不占缓存也不多拷一次。

\subsection{写路径}

\begin{txtcode}
fat32_write(node, offset, src, count)
+-- offset > size → 拒绝（不支持稀疏文件）
+-- 循环
    +-- ibmap(ip, pos / 簇大小, alloc=1)   不够就 fat_alloc() 接到链尾
    +-- 写满一个扇区 → bget_blank()；否则 bread()（读改写）
    +-- copy → bwrite()（写直达）→ brelse
    +-- 文件变长 → ip->size 更新, dirty = 1
\end{txtcode}

\begin{itemize}
  \item \code{fat\_alloc(prev)}：从 \code{free\_hint} 开始线性扫 FAT 找值为 0 的簇，标成链尾 \code{0x0FFFFFFF}，再把 \code{prev} 指向它。
  \item \code{fat\_set()} 同时写\textbf{所有 FAT 副本}（通常 2 份），保持一致，Windows 和 macOS 的 fsck 不会报 FAT 不一致。
  \item 写 FAT 项只改低 28 位，保留高 4 位（规范要求）。
  \item \textbf{FSInfo}：第一次分配簇之前把 FSInfo 里的``空闲簇数 / 下一个空闲簇''标成未知（\code{0xFFFFFFFF}），让其它系统下次挂载时自己重新计算。这样内核就不需要维护精确的空闲计数。
\end{itemize}

\subsection{创建 \code{fat32\_create()}（\code{open} 带 \code{O\_CREAT}）}

\begin{enumerate}
  \item \code{resolve(path, parent=1)} 找到父目录和最后一个分量。
  \item 文件已存在就直接打开；否则 \code{create\_entry()}：
    \begin{itemize}
      \item 名字能直接放进 8.3（\code{fit\_short()}，大写字母、合法字符、基名 $\leq$ 8、扩展名 $\leq$ 3；全小写的名字用 NT 大小写标志记录）→ 只写一个短目录项；
      \item 否则生成别名 \code{NAME$\sim$N.EXT}（\code{make\_alias()}，N 从 1 递增直到不冲突），再写 $\lceil$长度/13$\rceil$ 个 LFN 项 + 1 个短目录项，LFN 项带短名校验和。
      \item \code{find\_slots()} 在父目录里找足够的连续空位（\code{0x00} / \code{0xE5}），不够就给目录追加一个清零的新簇。
      \item 创建日期写固定的 2026-01-01（\code{0x5c21}），没有 RTC。
    \end{itemize}
  \item 新文件首簇为 0、大小为 0，第一次写入时才分配簇。
\end{enumerate}

\subsection{截断 \code{fat32\_truncate()}（\code{open} 带 \code{O\_TRUNC}）}

顺序是\textbf{先改目录项、再释放簇链}：首簇和大小清零 → \code{iupdate()} 写回 → \code{fat\_free\_chain()}。如果中途掉电，最坏情况是留下一条无人引用的簇链（fsck 能回收），而不会出现目录项指向已经释放的簇。

\subsection{列目录 \code{fat32\_readdir()}}

跳过隐藏项、卷标、\code{.} 和 \code{..}；类型取值和 minifs 一致（1 = 目录，2 = 文件），所以现有的 \code{ls} 不用改。\code{vfs\_dirent.name} 只有 31 字节，超过 30 个字符的长文件名在 \code{ls} 里会被截断显示（打开文件时用完整名字没问题）。

\subsection{不支持}

\code{mkdir}、\code{unlink}、\code{rename}、时间戳更新、稀疏文件、多分区（只认 MBR 里第一个 FAT32 分区）、FAT12/16、exFAT。

\section{VFS 的配套改动}

\file{include/fs/vfs.h}：

\begin{ccode}
#define VFS_O_ACCMODE 3u
#define VFS_O_CREATE 0x200u      /* xv6 fcntl.h values */
#define VFS_O_TRUNC  0x400u

struct vfs_backend_ops {
    ...
    /* Optional: create an empty regular file (O_CREATE). */
    int (*create)(const char *path, struct vfs_node *node);
    /* Optional: discard a file's contents (O_TRUNC). */
    int (*truncate)(struct vfs_node *node);
    /* Optional: the last reference to an open node is gone. */
    void (*close)(struct vfs_node *node);
};
\end{ccode}

\code{t113\_vfs\_open()} 的流程：

\begin{txtcode}
find_mount(path)                    最长前缀匹配
+-- lookup(path) 失败
|   +-- 带 O_CREATE 且后端有 create → create(path)
+-- 带 O_TRUNC（必须是 O_RDWR）→ truncate(node)
+-- 分配 open_files 槽位
t113_vfs_write()   成功后更新 file->node.size
t113_vfs_close()   引用计数归零 → ops->close(node)   fat32 在这里写回目录项
\end{txtcode}

三个新回调都是可选的：minifs、devfs、bootfs 不实现它们，行为不变（对它们带 \code{O\_CREATE} 会失败）。

用户态：\file{fcntl.h} 增加 \code{O\_CREATE} / \code{O\_CREAT}（0x200）和 \code{O\_TRUNC}（0x400），\code{sys\_open} 接受这两个标志。

\section{谁在用 \file{/boot}}

\subsection{XR829 固件}

驱动代码没有改：\code{need\_image()}（现在在 \file{xradio/fwio.c}）一直是

\begin{ccode}
struct vfs_file *file = t113_vfs_open(
    "/boot/fw_xr829.bin", VFS_O_RDONLY);
\end{ccode}

变的是 VFS 下面的后端。现在这个调用由 fat32 处理：在 FAT32 根目录里按名字查找，沿簇链从卡上读出。也就是说固件加载现在走的是真正的文件系统查找；只有 FAT32 挂载失败时才会退回 bootfs 的 DDR 视图。

\textbf{注}：SPL 里的 \code{load\_fat32\_firmware()} 仍然会把三份固件预读进 DDR，内核正常路径已经不依赖它，只给 bootfs 后备路径用。上板确认 FAT32 稳定后可以删掉这段，启动会快一点。

\subsection{tftp}

\file{src/net/tftp/tftp\_udp.c} 的 \code{tftp\_get()}：先把整个文件收进内存（最多 8 MiB），传输完成后才 \code{t113\_vfs\_open(local, O\_RDWR | O\_CREATE | O\_TRUNC)} 写入，关闭后再重新打开读回、比对 CRC32。传输中断不会在 \file{/boot} 留下写了一半的文件。用户命令 \code{tftp <服务器IP> <远端文件> [本地路径]}，本地路径默认 \file{/boot/<文件名>}。

\subsection{WiFi 配置}

supplicant（\file{src/net/wpa/supplicant.c} 的 \code{wpa\_supplicant\_connect()}）从 \file{/boot/wifi.txt} 读 \code{ssid=}、\code{psk=}（以及可选的 \code{ip=}、\code{netmask=}、\code{gateway=}）。密码只放在 FAT32 分区上，不进源码；仓库里对应的 \file{rpi\_bm/firmware/wifi.txt} 被 \code{.gitignore} 忽略。

\subsection{SPL 从 FAT32 加载内核}

miniBoot \file{spl.c} 新增 \code{load\_fat32\_kernel()}：SD 初始化后先在 FAT32 根目录找 \code{kernel\_minios.img}（最大 8 MiB），读到 \code{0x40200000}，打印大小和 CRC32；找到了就跳过 bootpack 里的 kernel 条目，找不到就按原来的 bootpack 流程走。配合 tftp，可以在板子上直接更新内核：

\begin{shcode}
tftp 192.168.0.100                        # 写入 /boot/kernel_minios.img
# 重启后 SPL 加载新内核
\end{shcode}

想回到 bootpack 里的内核，删掉 FAT32 分区上的 \code{kernel\_minios.img} 即可。

\section{FAT32 图解（参照 xv6 book 第 8 章）}

xv6 book 讲文件系统时，先给一张分层图，再给磁盘布局，然后逐层讲buffer cache、inode、目录、路径名，最后串起 \code{open/read/write}。本节按同样的顺序讲 miniOS 的 FAT32，每一层都和 xv6 的对应层并排比较。

\subsection{分层}

\begin{txtcode}
        xv6 文件系统（7 层）                     miniOS /boot（T113）
 +-------------------------------+      +--------------------------------------+
 | 文件描述符  file.c sysfile.c  |      | task->files[4] → struct vfs_file     | vfs.c
 +-------------------------------+      +--------------------------------------+
 | 路径名      namei / namex     |      | find_mount() 选后端（最长前缀）       | vfs.c
 +-------------------------------+      | resolve() 逐级 dir_find()            | fat32.c
 | 目录        dirlookup/dirlink |      +--------------------------------------+
 +-------------------------------+      | 目录项 32 B 短名 + LFN 长名           | fat32.c
 | inode       iget/ilock/bmap   |      | fat_inode（内存）+ ibmap() 走簇链     | fat32.c
 +-------------------------------+      +--------------------------------------+
 | 日志        begin_op/end_op   |      | （无日志：直写 + 写入顺序约定）        |
 +-------------------------------+      +--------------------------------------+
 | buffer cache  bread/bwrite    |      | bio.c  24 个 512 B 扇区缓冲，LRU      | bio.c
 +-------------------------------+      +--------------------------------------+
 | 磁盘        virtio_disk.c     |      | t113_sd_read/write  SMHC0 CMD17/CMD24 | sunxi-sd.c
 +-------------------------------+      +--------------------------------------+
\end{txtcode}

xv6 的 inode 层同时负责“文件数据在哪些块”（\code{addrs[]} 加间接块）和“文件元数据存在哪”（磁盘 inode 区）。FAT32 把这两件事拆开了：

\begin{itemize}
  \item 数据块的位置记在 FAT 表里，一簇接一簇连成链表；
  \item 文件大小和首簇号记在父目录的 32 字节目录项里。
\end{itemize}

所以 miniOS 的 \code{fat\_inode} 只是内存里的一个句柄，记住“我的目录项在磁盘哪个位置”，关闭文件时把首簇和大小写回去。

\subsection{磁盘布局}

整张卡和 FAT32 分区内部：

\begin{txtcode}
整张 SD 卡（512 B 扇区，LBA = 扇区号）
LBA 0          MBR：446 B 引导码 | 4 × 16 B 分区表 | 55 AA
               表项 1  type 0x0b/0x0c   → FAT32  → 挂到 /boot
               表项 2  type 0x83        → minifs → 挂到 /
LBA 16         Boot0 SPL（8 KiB 处）
LBA 256        bootpack（128 KiB 处）
part_lba       FAT32 分区开始 -------------+
                                           v
part_lba   +---------------------------------+ +
           | 扇区 0  引导扇区（BPB）          | |
           | 扇区 1  FSInfo                  | | 保留区：BPB_RsvdSecCnt 个扇区
           | 扇区 6  备份引导扇区             | | （常见 32）
           | ...                              | |
fat_lba    +---------------------------------+ + = part_lba + RsvdSecCnt
           | FAT #1                          |   sectors_per_fat 个扇区
           | 每个簇一项，4 字节，低 28 位有效 |   一个扇区放 128 项
           +---------------------------------+
           | FAT #2（镜像）                   |   fat_set() 两份都写
data_lba   +---------------------------------+ = fat_lba + fats × sectors_per_fat
           | 簇 2   根目录通常从这里开始       |   cluster_lba(c) = data_lba + (c - 2) × spc
           | 簇 3                            |
           | ...                              |
           | 簇 clusters+1                   |
           +---------------------------------+
\end{txtcode}

xv6 的 \code{mkfs} 布局是连续的几段，块号由超级块给出：

\begin{txtcode}
xv6 (BSIZE = 1024)
[ 0 boot | 1 super | 2.. log | inodestart.. inode 块 | bmapstart.. 位图 | 数据块 ... ]
\end{txtcode}

两者对照：

\begin{table}[H]
\centering
\small
\begin{tabularx}{\linewidth}{>{\raggedright\arraybackslash}X>{\raggedright\arraybackslash}X>{\raggedright\arraybackslash}X}
\toprule
xv6 & FAT32 & miniOS 里在哪 \\
\midrule
超级块 \code{struct superblock} & 引导扇区里的 BPB & \code{t113\_fat32\_mount\_vfs()} 解析到全局 \code{fs} \\
空闲块位图 & FAT 表项为 0 表示空闲 & \code{fat\_alloc()} 从 \code{free\_hint} 线性扫描 \\
磁盘 inode 区 & 没有，元数据在目录项里 & 目录项的 \code{ent\_lba/ent\_off} \\
\code{addrs[12]} + 1 个间接块 & FAT 链表，长度不限 & \code{ibmap()} \\
日志区 & 没有 & 写入顺序约定（见后文） \\
\bottomrule
\end{tabularx}
\end{table}

\code{t113\_fat32\_mount\_vfs()} 用到的 BPB 字段：

\begin{table}[H]
\centering
\small
\begin{tabularx}{\linewidth}{lll>{\raggedright\arraybackslash}X}
\toprule
偏移 & 长度 & 字段 & 校验 / 用途 \\
\midrule
11 & 2 & BytsPerSec & 必须等于 512（\code{BSIZE}） \\
13 & 1 & SecPerClus & \code{fs.spc}，非 0 \\
14 & 2 & RsvdSecCnt & \code{fat\_lba = part\_lba + RsvdSecCnt} \\
16 & 1 & NumFATs & \code{fs.fats}，非 0 \\
17 & 2 & RootEntCnt & 必须为 0，这是 FAT32 的标志 \\
32 & 4 & TotSec32 & 算数据簇数 \code{clusters} \\
36 & 4 & FATSz32 & \code{sectors\_per\_fat} \\
44 & 4 & RootClus & \code{root\_cluster}，必须是合法簇号 \\
48 & 2 & FSInfo & FSInfo 扇区号，0 表示没有 \\
510 & 2 & 签名 & 55 AA \\
\bottomrule
\end{tabularx}
\end{table}

簇数还会再被 FAT 表的容量截一次：\code{clusters ≤ sectors\_per\_fat × 128 - 2}。有些格式化工具算出来的 FAT 偏小，截一下能防止 \code{fat\_get()} 越过 FAT 表。

\subsection{FAT 表与簇链}

FAT 表就是一个大数组，下标是簇号，值是“下一个簇号”：

\begin{txtcode}
目录项：name = "fw_xr829.bin"  first = 5  size = 174624

FAT[]  下标:  2     3     4     5     6     7     8     9     10    11
       值:   EOC   EOC   0     6     7     9     0     10    EOC   0
                         空闲  |     |     |     空闲  |     |     空闲
                               +-→ 6 +-→ 7 +---------→ 9 +-→ 10 → EOC

文件第 k 个簇 = 从 first 出发沿链走 k 步
偏移 off 所在的簇序号 = off / (spc × 512)
扇区 = cluster_lba(簇号) + (off % 簇字节) / 512
\end{txtcode}

（簇号只是示意。）表项取值：

\begin{table}[H]
\centering
\small
\begin{tabularx}{\linewidth}{l>{\raggedright\arraybackslash}X}
\toprule
值（低 28 位） & 含义 \\
\midrule
\code{0} & 空闲 \\
\code{2 … clusters+1} & 下一个簇 \\
\code{≥ 0x0ffffff8} & 链尾 EOC，miniOS 写 \code{0x0fffffff} \\
\bottomrule
\end{tabularx}
\end{table}

\code{ibmap(ip, index, alloc)} 把“文件第 index 个簇”翻译成簇号，相当于 xv6 的\code{bmap()}。差别在于：xv6 用下标直接查 \code{addrs[]}，O(1)；FAT 要从头走链，O(n)。miniOS 在 \code{fat\_inode} 里放了一个游标 \code{cur\_index/cur\_cluster}，顺序读写时从上一次的位置接着走，所以整文件顺序读还是 O(n) 而不是 O(n²)：

\begin{txtcode}
ibmap(ip, index, alloc)
+-- ip->first == 0（空文件）
|   +-- alloc ? first = fat_alloc(0) : 返回 0
+-- 起点：cur_index ≤ index 时从游标出发，否则从 first 出发
+-- while i < index:
|   +-- next = fat_get(c)               bread(FAT 扇区)，取 4 字节
|   +-- next 不是合法簇
|   |   +-- 是 EOC 且 alloc → next = fat_alloc(c)   扩链
|   |   +-- 否则返回 0（越过文件末尾）
|   +-- c = next, i++
+-- 更新游标，返回 c

fat_alloc(prev)
+-- fsinfo_invalidate()                 第一次分配时把 FSInfo 计数写成“未知”
+-- 从 free_hint 起环形扫描 FAT，找值为 0 的项
+-- fat_set(c, EOC)                     所有 FAT 副本
+-- prev ? fat_set(prev, c)             挂到链尾
+-- free_hint = c + 1
\end{txtcode}

\subsection{目录项}

FAT32 的目录也是文件，内容是一串 32 字节目录项。一个 512 B 扇区放 16 项，一个簇放 \code{spc × 16} 项。短目录项（8.3）的格式：

\begin{txtcode}
偏移  长度  字段
 0     8   主名，空格填充      e[0] = 0x00 目录结束；0xE5 已删除；0x05 表示首字符是 0xE5
 8     3   扩展名
11     1   属性   0x01 只读 0x02 隐藏 0x04 系统 0x08 卷标 0x10 目录 0x20 归档；0x0F = LFN
12     1   NT 保留  bit3 主名全小写，bit4 扩展名全小写
13     1   创建时间（10 ms）
14     2   创建时间
16     2   创建日期              miniOS 新建文件写 2026-01-01
18     2   访问日期
20     2   首簇号高 16 位   +
22     2   修改时间         |
24     2   修改日期         |
26     2   首簇号低 16 位   + first = (hi << 16) | lo
28     4   文件大小（目录为 0）
\end{txtcode}

长文件名（LFN）放在短目录项前面，每项 13 个 UTF-16 字符，按逆序存放：

\begin{txtcode}
磁盘顺序 →
+---------------+---------------+------------------------+
| LFN  0x42     | LFN  0x01     | 短项 "MINIOS~1BIN"     |
| 字符 14..17   | 字符 1..13    | attr 0x20 first size   |
| attr 0x0F     | attr 0x0F     |                        |
| chksum = S    | chksum = S    | S = short_checksum(名) |
+---------------+---------------+------------------------+
  0x40 标记“最后一片”（也就是磁盘上第一项）

LFN 项内 13 个字符的字节位置：lfn_pos[] = {1,3,5,7,9, 14,16,18,20,22,24, 28,30}
校验和：sum = ((sum & 1) << 7) + (sum >> 1) + name11[i]，i = 0..10
\end{txtcode}

读目录时，\code{dir\_next()} 是一个小状态机：

\begin{txtcode}
dir_next(it, d)
+-- 本簇的项读完 → it->cluster = fat_get(it->cluster)
+-- bread(cluster_lba(c) + index / 16)，拷出 32 字节，brelse
+-- e[0] == 0x00   → 目录结束，返回 0
+-- e[0] == 0xE5   → 已删除，丢弃攒了一半的长名
+-- attr == 0x0F   → lfn_piece()：0x40 项清空缓冲并记下 chksum，按 ord 填 13 个字符
+-- attr & VOLUME  → 卷标，跳过
+-- 短项
    +-- 长名完整且 chksum 匹配 → d->name = 长名
    +-- 否则 short_name()     → "FW_XR829.BIN"，按 NT 位转小写
    +-- d->first / size / lba / off，返回 1
\end{txtcode}

新建文件时，名字能精确表示成 8.3 就只写一个短项。例如 \code{fw\_xr829.bin} 存成\code{FW\_XR829BIN}，NT 字节写 0x18 表示主名和扩展名都是小写。表示不了时（太长、多个点、大小写混合等）才生成 \code{NAME\textasciitilde{}N.EXT} 别名并写 LFN 项：

\begin{txtcode}[basicstyle=\ttfamily\scriptsize]
create_entry(dir_cluster, "minios-kernel.bin")
+-- long_char_ok() 逐字符检查         不允许 \ / : * ? " < > |
+-- fit_short() 失败（主名 13 字符 > 8）
+-- make_alias()                       "MINIOS~1BIN"，~1 已存在就试 ~2 ...
+-- pieces = ceil(17 / 13) = 2
+-- find_slots(dir, 3)                 找 3 个连续空项；目录满了就 fat_alloc + zero_cluster 扩一簇
+-- write_entry() × 2                  LFN 0x42、0x01
+-- write_entry()                      短项：attr ARCHIVE，first = 0，size = 0
\end{txtcode}

xv6 的目录项是 16 字节的 \code{struct dirent \{ ushort inum; char name[14]; \}}，\code{inum} 指向 inode，长度和数据都在 inode 里。FAT 的目录项本身就是 inode。

\subsection{内存中的数据结构}

\begin{txtcode}
用户 fd / 内核指针
      |
      v
task->files[i] --→ struct vfs_file              open_files[32]，vfs.c
                   +-- references / flags / offset
                   +-- ops --→ fat32_ops
                   +-- node { id = first, size, type,
                              private_data } --+
                                               v
                             fs.inode[16]：struct fat_inode      fat32.c
                             +-- ref                    同一目录项只占一个槽（iget 去重）
                             +-- dir / first / size
                             +-- ent_lba / ent_off --→ 磁盘上短目录项的位置（根目录为 0）
                             +-- dirty                  first/size 改过，close 时 iupdate()
                             +-- cur_index / cur_cluster  簇链游标
                                               | fat_get / bread / bwrite
                                               v
                             bcache：struct buf[24]               bio.c
                             head <-> 最近释放 <-> ... <-> 最久未用 <-> head
                             +-- bget：正向找命中（valid && blockno 相同）
                             |         反向找 refcnt == 0 的回收
                             +-- bread：未命中时 t113_sd_read(blockno, 1)
                             +-- bwrite：立刻 t113_sd_write（直写）
                             +-- brelse：refcnt 归 0 时移到 head.next
                                               |
                                               v
                             sunxi-sd.c：SMHC0，PF0–PF5，1-bit，12 MHz
                             CMD17 单块读 / CMD24 单块写，CPU 搬 FIFO，sd_lock 互斥
\end{txtcode}

与 xv6 的 \code{bio.c} 相比：

\begin{table}[H]
\centering
\small
\begin{tabularx}{\linewidth}{l>{\raggedright\arraybackslash}X>{\raggedright\arraybackslash}X}
\toprule
方面 & xv6 \code{bio.c} & miniOS \code{bio.c} \\
\midrule
块大小 & 1024 B（\code{BSIZE}） & 512 B，正好一个 SD 扇区 \\
缓冲数 & \code{NBUF = MAXOPBLOCKS × 3 = 30} & 24 \\
每块锁 & 每个 \code{buf} 一把 sleeplock，\code{bread} 返回时已上锁 & 没有每块锁；FAT32 层的 \code{fs.lock} 一把互斥锁串行所有操作 \\
写 & \code{bwrite} 由日志层决定何时写，\code{log\_write} 钉住缓冲 & 直写，\code{bwrite} 立刻写卡；失败时把缓冲标成无效，下次重读 \\
大块读 & 总是经过缓存 & \code{fat32\_read} 对齐的整扇区直接 \code{t113\_sd\_read} 进调用者缓冲，绕过缓存 \\
设备 & virtio，中断完成，进程 sleep & SMHC0 轮询，忙等 \\
\bottomrule
\end{tabularx}
\end{table}

绕过缓存是安全的，因为缓存是直写：卡上的内容永远不比缓存旧。

\subsection{调用链：打开并读一个文件}

XR829 驱动装固件时的路径（\code{fwio.c} 的 \code{need\_image()}）：

\begin{txtcode}
need_image("/boot/fw_xr829.bin")
+-- t113_vfs_open(path, VFS_O_RDONLY)                         vfs.c
|   +-- mutex_lock(&vfs_lock)
|   +-- find_mount(path)            "/" 和 "/boot" 都匹配，取更长的 "/boot"
|   +-- fat32_lookup(path, &node)                             fat32.c
|   |   +-- mutex_lock(&fs.lock)
|   |   +-- resolve(path, &d, parent = 0)
|   |   |   +-- mount_relative()    "/boot/fw_xr829.bin" → "fw_xr829.bin"
|   |   |   +-- d = 根目录（first = root_cluster，lba = 0）
|   |   |   +-- dir_find(root_cluster, "fw_xr829.bin")
|   |   |       +-- while dir_next() > 0: name_equal()        不区分大小写
|   |   |           +-- bread → 命中或 t113_sd_read(lba, 1)
|   |   +-- iget(dir = 0, first, size, lba, off)              按目录项位置去重
|   |   +-- fill_node()             node.private_data = fat_inode
|   +-- 取一个空闲 open_files[]，offset = 0
+-- t113_vfs_file_size(file)        174624
+-- t113_alloc_pages(43)
+-- t113_vfs_read(file, data, 174624)
|   +-- fat32_read(&node, 0, data, 174624)
|       +-- while done < count
|           +-- c = ibmap(ip, pos / 簇字节, 0)    游标让顺序读每簇只走一步
|           +-- 扇区对齐、剩余 ≥ 512、dst 4 字节对齐
|           |   +-- t113_sd_read(lba, n, dst + done)   一次读到本簇末尾
|           +-- 否则 bread(lba) → copy → brelse        首尾不对齐的零头
+-- t113_vfs_close(file)
    +-- fat32_close()               未修改，只做 ref--
\end{txtcode}

\subsection{调用链：创建并写一个文件}

\code{tftp} 把下载的文件写进 \file{/boot} 时（\code{src/net/tftp/tftp\_udp.c} 的\code{write\_file()}）：

\begin{txtcode}
t113_vfs_open("/boot/minios-kernel.bin", RDWR | CREATE | TRUNC)
+-- find_mount → fat32_ops
+-- fat32_lookup() 失败（文件还不存在）
+-- fat32_create(path, &node)
|   +-- resolve(path, &d, parent = 1, &name, &len)    停在父目录
|   +-- dir_find() 再确认不存在
|   +-- create_entry()                                LFN + 短项，first = 0
|   +-- iget(0, 0, 0, lba, off)
+-- fat32_truncate(&node)          新文件也走一遍
    +-- first = 0, size = 0, 清游标
    +-- iupdate(ip)                先改目录项...
    +-- fat_free_chain(旧 first)   ...再释放旧簇链

t113_vfs_write(f, data, 4096) × N
+-- fat32_write(&node, offset, src, count)
    +-- offset > size → -1          不支持稀疏文件
    +-- while done < count
        +-- c = ibmap(ip, pos / 簇字节, alloc = 1)
        |   +-- 走到链尾 → fat_alloc(c)
        |       +-- fat_set(new, EOC)   FAT1、FAT2 各 bread + bwrite
        |       +-- fat_set(c, new)
        +-- 整扇区覆盖 → bget_blank(lba)   不必先读
        |   部分扇区   → bread(lba)
        +-- copy → bwrite → brelse       数据扇区立刻落盘
        +-- size 变大 → dirty = 1

t113_vfs_close(f)
+-- fat32_close()
    +-- dirty → iupdate()            目录项写回 first 和 size
    +-- ref--
\end{txtcode}

\subsection{没有日志时的写入顺序}

xv6 的每个文件系统调用包在 \code{begin\_op()/end\_op()} 里。日志提交前，所有修改只存在于缓存和日志区，崩溃后要么全部重放，要么全部丢弃。miniOS 的 FAT32没有日志，靠写入顺序把崩溃的后果限制在“丢空间，不丢一致性”：

\begin{table}[H]
\centering
\small
\begin{tabularx}{\linewidth}{l>{\raggedright\arraybackslash}X>{\raggedright\arraybackslash}X}
\toprule
操作 & 顺序 & 中途断电的后果 \\
\midrule
追加写 & 先 \code{fat\_alloc} 写 FAT，再写数据扇区，\code{close} 时最后写目录项 & 目录项还是旧的 first/size，新簇成为没人引用的“丢失簇”；\code{fsck} 能回收 \\
截断 & 先 \code{iupdate} 把目录项清零，再 \code{fat\_free\_chain} & 旧链没释放完，同样只是丢失簇 \\
新建 & 先写 LFN 项，再写短项 & 只有 LFN 没有短项时，校验和对不上，读目录会忽略这几项 \\
FSInfo & 第一次分配时把空闲计数写成 \code{0xffffffff} & 其他系统会重新数空闲簇，不会相信过期的计数 \\
\bottomrule
\end{tabularx}
\end{table}

FAT 的两份副本是先后写的，中途断电时两份可能不一致。Windows 和 \code{fsck.vfat}默认以第一份为准。

\subsection{T113 实现上的区别}

\begin{itemize}
  \item 块设备是 SD 卡本身：\code{bio.c} 的块号就是整卡 LBA，和 xv6 的“文件系统内块号”不同。FAT32 挂载时从 MBR 找第一个类型为 0x0b/0x0c/0x1b/0x1c 的分区，所有地址都加上 \code{part\_lba}。
  \item \code{sunxi-sd.c} 和 SPL 共用：SMHC0 走 PF0–PF5，1-bit 总线，12 MHz，每次一块（CMD17/CMD24），CPU 轮询 FIFO。内核里加了 \code{sd\_lock} 互斥锁，根文件系统（\code{rootdev.c}）和 \file{/boot} 共享这一个控制器。
  \item 一把锁：\code{fs.lock} 串行 FAT32 的所有操作，外面还有 VFS 的 \code{vfs\_lock}。没有xv6 那种每个 inode、每个缓冲一把 sleeplock 的细粒度并发。
  \item 不支持：mkdir、unlink、rename、时间戳（固定写 2026-01-01）、稀疏写。目录只能读和在已有目录里新建文件。
  \item 失败可回退：\code{t113\_fat32\_mount\_vfs()} 失败时，\code{kernel\_boot()} 改挂\code{bootfs}，用 SPL 预加载到 DDR 的固件镜像（见 第~\ref{chap:rootfs} 章）。
\end{itemize}

\section{一致性与限制}

\begin{itemize}
  \item \textbf{没有日志}。写直达保证每次 \code{write} 返回时数据已经在卡上，但目录项里的大小和首簇只在最后一次 \code{close} 时写回。如果写到一半掉电，新分配的簇在 FAT 里已经占用，目录项却还是旧的大小，表现为文件内容没变、分区里多出一些丢失的簇，fsck（\code{fsck.fat} 或 macOS 磁盘工具）能回收。
  \item \textbf{锁}：整个 fat32 一把互斥锁（类似早期 xv6 的大锁），每个 VFS 回调从头锁到尾。互斥锁不可重入，fat32 内部函数之间不会重复加锁。
  \item \textbf{inode 表 16 项}：同时打开的不同 FAT32 文件最多 16 个（VFS 的 \code{open\_files} 也只有 32 个槽）。
  \item \textbf{SD 访问}：单块 CMD17/CMD24 轮询，没有 DMA、没有多块传输，大文件写入速度受此限制。
\end{itemize}

\section{改动前后对照}

\begin{table}[H]
\centering
\begin{tabular}{lp{4.6cm}p{5.2cm}}
\toprule
 & 改动前（bootfs） & 改动后（fat32） \\
\midrule
谁解析 FAT32 & SPL & 内核（SPL 也保留自己的只读解析器） \\
数据来源 & 启动时 DDR 快照 & SD 卡，运行期实时读写 \\
查找 & 三个硬编码路径 & 任意路径、长文件名、子目录 \\
\code{ls /boot} & 空 & 列出分区内容 \\
写入 / 创建 / 截断 & $\times$ & $\checkmark$ \\
缓存 & 无（数据本来就在内存） & bio 24 块 LRU，写直达 \\
失败时 & PANIC & 退回 bootfs，再失败才 PANIC \\
与 Linux 对照 & firmware loader 读 initramfs 里的文件 & vfat 驱动挂载 \file{/boot} \\
\bottomrule
\end{tabular}
\end{table}

\section{测试}

主机测试（\code{gcc -fsanitize=address,undefined}，链接内核原样的 \file{vfs.c}、\file{bio.c}、\file{fat32.c}，SD 读写换成对镜像文件的读写）：

\begin{itemize}
  \item \code{mkfs.fat -F 32} 生成分区（每簇扇区数作为参数，测过不同簇大小）；
  \item 读 mtools 写入的长文件名文件和子目录里的文件，逐字节比对；
  \item \code{O\_CREAT|O\_TRUNC} 新建文件，用 512 字节和 300 字节交替写入（覆盖对齐和不对齐、跨扇区、跨簇），读回比对；
  \item 截断已有的 \code{fw\_xr829.bin} 并重写；
  \item 在子目录里连续创建 40 个长文件名文件，触发目录追加新簇；
  \item 用 \code{fsck.fat -n} 检查分区无错误，再用 \code{mcopy} 读回内核写的文件比对。
\end{itemize}

上板预期串口输出：

\begin{txtcode}
SPL: FAT32 kernel kernel_minios.img bytes=0x... crc32=0x...   # FAT32 上有内核时
SPL: skip bootpack kernel (FAT32 kernel used)
...
fat32: /boot mounted rw, partition lba=0x00002000 clusters=0x... cluster-bytes=0x...
mmc1:0001:1: loaded /boot/boot_xr829.bin bytes 2184
...
$ ls /boot
$ tftp 192.168.0.100 test.bin
tftp: received ... bytes in ... ms, crc32 0x...
tftp: saved /boot/test.bin
\end{txtcode}

如果看到 \code{fat32: mount failed, using SPL handoff images}，说明 MBR 里没有 FAT32 分区或 BPB 校验失败，此时 \file{/boot} 退回老的只读视图。
